\section{ KẾT LUẬN}


\subsection{Kết quả đạt được}
Thông qua quá trình nghiên cứu và thực hiện đồ án, nhóm đã đạt được các kết quả sau:
\begin{itemize}
    \item \textbf{Thiết kế hệ thống:} Hoàn thành việc thiết kế và mô phỏng một RISC CPU đơn giản dựa trên ngôn ngữ mô tả phần cứng Verilog. 
    \item \textbf{Kiến trúc phân khối:} Học tập các cấu trúc khối chức năng và khối điều khiển của một RISC CPU.
    \item \textbf{Năng lực chuyên môn:} Quen với việc sử dụng ngôn ngữ Verilog để đặc tả các sơ đồ khối, đồng thời học cách tạo các testbench một cách có chủ đích và qua đó
    nhóm còn học được cách viết một báo cáo gồm đầy đủ các nội dung.
\end{itemize}

\subsection{Khó khăn và Thách thức}
Trong quá trình triển khai, nhóm đã đối mặt và giải quyết các thách thức kỹ thuật sau:
\begin{itemize}
    \item \textbf{Bước đầu làm quen:} Choáng ngợp trước mô hình phức tạp, có quá nhiều khối chức năng và các đường truyền dễ lẫn lộn.
    \item \textbf{Làm việc nhóm: } Việc làm việc nhóm với một đề tài tương đối lớn gây nhiều khó khăn về mặt phân chia công việc và xác định deadline.
\end{itemize}

\subsection{Hướng phát triển và Mở rộng}
Nhằm nâng cao hiệu suất và tính ứng dụng của kiến trúc, nhóm đề xuất các hướng phát triển tiếp theo:
\begin{itemize}
    \item \textbf{Nâng cấp tập lệnh (ISA):} Mở rộng khả năng xử lý bằng cách tích hợp các phép toán phức tạp (nhân/chia), xử lý số dấu phẩy động và các chế độ địa chỉ gián tiếp.
    \item \textbf{Tối ưu hóa kiến trúc Pipeline:} Triển khai cấu trúc đường ống nhiều giai đoạn (IF, ID, EX, MEM, WB) nhằm tăng cường thông lượng (throughput) và khai thác tối đa tài nguyên phần cứng.
    \item \textbf{Triển khai phần cứng thực tế:} Thực hiện tổng hợp và nạp thiết kế lên bo mạch FPGA (như Arty Z7) để đánh giá các thông số về tài nguyên thực tế và tần số hoạt động tối đa. Đồng thời phát triển trình hợp dịch (Assembler) riêng để hỗ trợ lập trình mức cao cho CPU.
\end{itemize}

\subsection{Bảng phân chia công việc}
\begin{table}[H]
\centering
\caption{Bảng phân chia công việc thực hiện đồ án}
\renewcommand{\arraystretch}{1.5}
% Sử dụng >{\raggedright\arraybackslash} để căn lề trái trong ô, giúp hết Underfull
\begin{tabularx}{\textwidth}{|c|>{\raggedright\arraybackslash}X|c|>{\raggedright\arraybackslash}p{6cm}|c|}
\hline
\textbf{STT} & \textbf{Họ và tên} & \textbf{MSSV} & \textbf{Nhiệm vụ} & \textbf{Hoàn thành} \\ \hline
1 & PHẠM ĐỖ HỒNG PHÚC & 2510359 & Thiết kế \& code Verilog khối ALU và Accumulator Register. Viết testbench tương ứng.Tổng hợp và ghi báo cáo latex. & 100\% \\ \hline
2 & TRẦN TIẾN THÀNH & 2413182 &  Thiết kế \& code Verilog khối Controller. Tích hợp Top-level và viết testbench toàn hệ thống. & 100\% \\ \hline
3 & VÕ LÊ PHÚC THỊNH & 2510382 & Thiết kế \& code Verilog khối Memory, PC, Address Mux và IR. Viết testbench tương ứng. & 100\% \\ \hline
\end{tabularx}
\end{table}